% !TEX root =  Lecture13.tex


%%%%%%%%%%%%%%%%%%%%%%%%%%%%%%%%%%%%
% Recap/Agenda
%%%%%%%%%%%%%%%%%%%%%%%%%%%%%%%%%%%%
\begin{frame}
\frametitle{Module 3: Regression Inference}
    \begin{itemize}
        \item \hl{Previously:} inference about \emph{group means}: ANOVA, and the two roads for the $F$-statistic. That closed Module 2.
        \item \hl{This time:} a new statistic, the \hl{regression slope} $\hat\beta_1$. We can fit a line (Lectures 3--7); now we ask whether its tilt is \emph{real}.
          \begin{itemize}
              \item the same two roads: shuffle the pairing, or a $t$ formula;
              \item a bootstrap confidence interval for the slope;
              \item the conditions (\hl{LINE}) and how to check them;
              \item the intercept, and testing a slope against a value other than zero.
          \end{itemize}
        \item \hl{Reading:} IMS Chapter 24.
        \item \hl{Homework:} the assigned reading.
    \end{itemize}
\end{frame}


%%%%%%%%%%%%%%%%%%%%%%%%%%%%%%%%%%%%
% Learning objectives:
%%%%%%%%%%%%%%%%%%%%%%%%%%%%%%%%%%%%
\begin{frame}
    \frametitle{Lecture Objectives}
    \goals{%
    \begin{itemize}
    \item test $H_0: \beta_1 = 0$ by \emph{shuffling the pairing} and with $t_{n-2}$, and say why they agree;
    \item build a slope CI by \emph{bootstrapping the pairs} and with $\hat\beta_1 \pm t^\ast\,\mathrm{SE}$;
    \item check the \hl{LINE} conditions on residual and Q--Q plots, and know which road survives failure;
    \item read an \texttt{lm()}/\texttt{tidy()} table, and interpret the intercept and a test of $\beta_1 = c$.
    \end{itemize}}
    \objectives{%
    \begin{itemize}
    \item \textbf{(M3, LO1)} Randomization Tests and Bootstrap Confidence Intervals in Regression
    \item \textbf{(M3, LO2)} Model Checking and Remedial Measures in Linear Regression
    \end{itemize}}
\end{frame}
